\begin{proposition}[Reflection reciprocity gives no second bridge condition]
\label{prop:bridge-reflection-reciprocity}
Let $p\ge5$ be prime, and let $u,w$ be two $p$-integral solutions of the
Ap\'ery recurrence whose reductions have nonzero Casoratian.  Put
\[
 \mathbf v_n=(n!^3u_n,n!^3w_n)\in\mathbf F_p^2
 \qquad(0\le n<p).
\]
Extend the gap-polynomial convention by $N_0=0$.  For
$0\le x<x+h<p$, one has
\begin{equation}
 \mathbf v_{x+h}
 =N_h(x)\mathbf v_{x+1}
  -(x+1)^6N_{h-1}(x+1)\mathbf v_x.
\label{eq:bridge-return-expansion}
\end{equation}
Consequently, if $N_h(x)=0$, then
\begin{equation}
 \mathbf v_{x+h}=\mu_h(x)\mathbf v_x,
 \qquad
 \mu_h(x)=-(x+1)^6N_{h-1}(x+1)\ne0.
\label{eq:bridge-return-multiplier}
\end{equation}

Define
\[
 \rho_n=(-1)^{n+1}(n!)^{-6}\in\mathbf F_p^\times.
\]
The pointwise reflection theorem gives
\begin{equation}
 \mathbf v_{p-1-n}=\rho_n\mathbf v_n.
\label{eq:bridge-row-reflection}
\end{equation}
If $x^\vee=p-1-x-h$ is the start of the reflected interval, then
$N_h(x^\vee)=0$ and
\begin{equation}
 \boxed{
 \mu_h(x)\mu_h(x^\vee)
 =\frac{\rho_x}{\rho_{x+h}}
 =(-1)^h\left(\prod_{j=1}^h(x+j)\right)^6.}
\label{eq:bridge-multiplier-reciprocity}
\end{equation}
In particular, multiplying a bridge transfer by its reflected bridge
transfer produces a prescribed unit, not an additional vanishing
condition.
\end{proposition}

\begin{proof}
After the factorial renormalization, each coordinate of $\mathbf v_n$
satisfies
\[
 \mathbf v_{n+1}=P(n)\mathbf v_n-n^6\mathbf v_{n-1}.
\]
Equation~\eqref{eq:bridge-return-expansion} follows by induction from this
recurrence and the defining recurrence for $N_h$.  The cases $h=1,2$ are
immediate from $N_0=0$, $N_1=1$, and $N_2(x)=P(x+1)$; the induction step
uses
\[
 N_{h+1}(x)=P(x+h)N_h(x)-(x+h)^6N_{h-1}(x)
\]
for the coefficient of $\mathbf v_{x+1}$, and the same identity with
$x$ replaced by $x+1$ for the coefficient of $\mathbf v_x$.
If $N_h(x)=0$, this proves~\eqref{eq:bridge-return-multiplier}.
The multiplier is nonzero because the Casoratian of the two solutions is
nonzero throughout the nonwrapping range, so the row $\mathbf v_{x+h}$
cannot vanish.

Proposition~\ref{prop:reflection} gives
$u_{p-1-n}=u_n$ and $w_{p-1-n}=w_n$.  Wilson's theorem gives
\[
 (p-1-n)!\,n!\equiv(-1)^{n+1}\pmod p,
\]
which proves~\eqref{eq:bridge-row-reflection} after cubing and comparing
with the normalization of $\mathbf v_n$.

Write $y=x+h$.  Reflection of
$\mathbf v_y=\mu_h(x)\mathbf v_x$ gives
\[
 \mathbf v_{p-1-y}=\rho_y\mu_h(x)\mathbf v_x,
 \qquad
 \mathbf v_{p-1-x}=\rho_x\mathbf v_x.
\]
The reflected interval runs from $x^\vee=p-1-y$ to
$x^\vee+h=p-1-x$, so its return multiplier satisfies
\[
 \rho_x\mathbf v_x
 =\mu_h(x^\vee)\rho_y\mu_h(x)\mathbf v_x.
\]
Cancelling the nonzero row and the units proves the first equality in
\eqref{eq:bridge-multiplier-reciprocity}.  The second follows directly
from the displayed formula for~$\rho_n$.
\end{proof}

\begin{remark}[Full-cycle multiplication is tautological]
\label{rem:bridge-full-cycle-tautology}
For two separated reflection orbits of short return chains, the bridge from
the first chain to the second and the bridge from the reflected second chain
to the reflected first chain form a reflected pair.  Formula
\eqref{eq:bridge-multiplier-reciprocity} shows that their multipliers
cancel up to the determinant unit already forced by the factorial gauge.
The connections pairing either chain with its reflected copy are themselves
instances of~\eqref{eq:bridge-row-reflection}.  If one writes a full-period
factorization, the remaining step crosses the singular index~$0$ and is
rank one; it cannot be inverted to manufacture another condition.  Direct
substitution therefore recovers the pointwise reflection identity.  It
neither removes the bridge length nor supplies a new zero congruence or a
new Sylvester-nullity factor.

Retaining the short span at both ends still yields the aligned determinant
formed from three polynomials in the long-bridge reduction; that
legitimate construction lowers the height for one fixed bridge label but
continues to depend on that label.  Reflection reciprocity offers only an
absolute-factor duplication and therefore cannot supply the missing
power-of-$X$ saving in the far-bridge sum.
\end{remark}
